% ============================================================
% Punkt 6: Źródła
% Projekt układu sumatora pełnego (Full Adder)
% z minimalizacją bramek logicznych
% ------------------------------------------------------------
% Plik przeznaczony do włączenia w dokumencie głównym Overleaf:
%   % ============================================================
% Punkt 6: Źródła
% Projekt układu sumatora pełnego (Full Adder)
% z minimalizacją bramek logicznych
% ------------------------------------------------------------
% Plik przeznaczony do włączenia w dokumencie głównym Overleaf:
%   % ============================================================
% Punkt 6: Źródła
% Projekt układu sumatora pełnego (Full Adder)
% z minimalizacją bramek logicznych
% ------------------------------------------------------------
% Plik przeznaczony do włączenia w dokumencie głównym Overleaf:
%   % ============================================================
% Punkt 6: Źródła
% Projekt układu sumatora pełnego (Full Adder)
% z minimalizacją bramek logicznych
% ------------------------------------------------------------
% Plik przeznaczony do włączenia w dokumencie głównym Overleaf:
%   \input{06_zrodla}
% ============================================================

\section{Źródła}

\begin{enumerate}
    \item Sasiadek J., \textit{Układy cyfrowe}, Wydawnictwa Naukowo-Techniczne, Warszawa 2018.
    \item Łuba T., \textit{Układy Logiczne w Zadaniach}, Oficyna Wydawnicza Politechniki Warszawskiej, Warszawa 2001.
    \item Gilmore C., \textit{Materiały dydaktyczne: Realizacja logicznych układów kombinacyjnych}, Instytut Fizyki Politechniki Łódzkiej.
    \item Wilkinson B., \textit{Układy cyfrowe}, Wydawnictwa Komunikacji i Łączności, Warszawa 2003.
    \item Górecki P., \textit{Układy cyfrowe, pierwsze kroki}, Wydawnictwo BTC, Warszawa 2004.
    \item Grobelna I., \textit{Układy cyfrowe -- laboratorium. Projektowanie układów kombinacyjnych na przykładzie komparatora}, Uniwersytet Zielonogórski.
    \item Toczek W., \textit{Układy cyfrowe -- projektowanie i symulacja}, Wydawnictwo Politechniki Łódzkiej, Łódź 2020.
\end{enumerate}

% ============================================================

\section{Źródła}

\begin{enumerate}
    \item Sasiadek J., \textit{Układy cyfrowe}, Wydawnictwa Naukowo-Techniczne, Warszawa 2018.
    \item Łuba T., \textit{Układy Logiczne w Zadaniach}, Oficyna Wydawnicza Politechniki Warszawskiej, Warszawa 2001.
    \item Gilmore C., \textit{Materiały dydaktyczne: Realizacja logicznych układów kombinacyjnych}, Instytut Fizyki Politechniki Łódzkiej.
    \item Wilkinson B., \textit{Układy cyfrowe}, Wydawnictwa Komunikacji i Łączności, Warszawa 2003.
    \item Górecki P., \textit{Układy cyfrowe, pierwsze kroki}, Wydawnictwo BTC, Warszawa 2004.
    \item Grobelna I., \textit{Układy cyfrowe -- laboratorium. Projektowanie układów kombinacyjnych na przykładzie komparatora}, Uniwersytet Zielonogórski.
    \item Toczek W., \textit{Układy cyfrowe -- projektowanie i symulacja}, Wydawnictwo Politechniki Łódzkiej, Łódź 2020.
\end{enumerate}

% ============================================================

\section{Źródła}

\begin{enumerate}
    \item Sasiadek J., \textit{Układy cyfrowe}, Wydawnictwa Naukowo-Techniczne, Warszawa 2018.
    \item Łuba T., \textit{Układy Logiczne w Zadaniach}, Oficyna Wydawnicza Politechniki Warszawskiej, Warszawa 2001.
    \item Gilmore C., \textit{Materiały dydaktyczne: Realizacja logicznych układów kombinacyjnych}, Instytut Fizyki Politechniki Łódzkiej.
    \item Wilkinson B., \textit{Układy cyfrowe}, Wydawnictwa Komunikacji i Łączności, Warszawa 2003.
    \item Górecki P., \textit{Układy cyfrowe, pierwsze kroki}, Wydawnictwo BTC, Warszawa 2004.
    \item Grobelna I., \textit{Układy cyfrowe -- laboratorium. Projektowanie układów kombinacyjnych na przykładzie komparatora}, Uniwersytet Zielonogórski.
    \item Toczek W., \textit{Układy cyfrowe -- projektowanie i symulacja}, Wydawnictwo Politechniki Łódzkiej, Łódź 2020.
\end{enumerate}

% ============================================================

\section{Źródła}

\begin{enumerate}
    \item Sasiadek J., \textit{Układy cyfrowe}, Wydawnictwa Naukowo-Techniczne, Warszawa 2018.
    \item Łuba T., \textit{Układy Logiczne w Zadaniach}, Oficyna Wydawnicza Politechniki Warszawskiej, Warszawa 2001.
    \item Gilmore C., \textit{Materiały dydaktyczne: Realizacja logicznych układów kombinacyjnych}, Instytut Fizyki Politechniki Łódzkiej.
    \item Wilkinson B., \textit{Układy cyfrowe}, Wydawnictwa Komunikacji i Łączności, Warszawa 2003.
    \item Górecki P., \textit{Układy cyfrowe, pierwsze kroki}, Wydawnictwo BTC, Warszawa 2004.
    \item Grobelna I., \textit{Układy cyfrowe -- laboratorium. Projektowanie układów kombinacyjnych na przykładzie komparatora}, Uniwersytet Zielonogórski.
    \item Toczek W., \textit{Układy cyfrowe -- projektowanie i symulacja}, Wydawnictwo Politechniki Łódzkiej, Łódź 2020.
\end{enumerate}
